\bookchapter{Bab 10: Pembayaran mobile dan layanan kehidupan}

\chapterintro{1. Bagaimana Cara Menghubungkan Kartu Bank?}

\booksection{Apakah Anda menggunakan Alipay atau WeChat, pertama kali menghubungkan kartu harus melalui langkah-langkah ini:}

Buka aplikasi, cari “Saya” atau “Dompet”. Ketuk “Tambah Kartu Bank”.

\begin{enumerate}
  \item Masukkan nomor kartu. Sekarang sebagian besar aplikasi mendukung foto langsung untuk mengenali nomor kartu, tidak perlu mengetik satu per satu.
  \item Bank akan mengirim SMS kode verifikasi ke nomor ponsel yang Anda daftarkan. Masukkan, lalu atur kata sandi pembayaran 6 digit.
  \item Kata sandi pembayaran ini, kata sandi kartu bank, dan kata sandi layar kunci Anda, sebaiknya tiga kata sandi yang berbeda, jangan disamakan demi kemudahan.
  \item Ada tip kecil: jika khawatir soal keamanan, bisa buka rekening elektronik tipe II atau III di aplikasi bank, khusus untuk menghubungkan pembayaran ponsel. Kartu ini hanya diisi uang belanja harian, meski terjadi masalah, kerugian juga terbatas.
\end{enumerate}

2. Kode Pembayaran Adalah “Siapa yang Melihat, Siapa yang Mengambil Uang”

\booksection{Banyak orang bingung membedakan kode pembayaran dan kode penerimaan, perbedaan ini sangat penting.}

Kode Pembayaran

\begin{itemize}
  \item \textbf{: Barcode + QR code yang Anda tunjukkan untuk dipindai kasir. Kode ini setara dengan "dompet elektronik" Anda, siapa yang memindainya bisa menarik uang dari rekening Anda. Jadi}jangan pernah screenshot dan kirim ke siapa pun\textbf{. Kode ini otomatis diperbarui setiap menit, justru untuk mencegah penyalahgunaan.}Kode Penerimaan
  \item \textbf{: Orang lain memindai kode ini untuk transfer ke Anda. Kode ini bisa di-screenshot disimpan, bisa dikirim ke orang lain, aman.}Pahami perbedaan ini, bisa menghindari banyak jebakan penipuan. Siapa pun yang meminta Anda "screenshot kode pembayaran dan kirimkan", seratus persen adalah penipu.
\end{itemize}

3. Pengaturan Keamanan Ini, 5 Menit Bisa Menghemat Banyak Kekhawatiran

\booksection{Pembayaran Sidik Jari/Wajah}

\begin{itemize}
  \item \textbf{: Di pengaturan pembayaran Alipay dan WeChat, aktifkan pembayaran biometrik. Sehingga setiap pembayaran perlu tekan sidik jari atau pindai wajah, meski ponsel hilang, orang yang menemukan tidak bisa membelanjakan uang Anda.}Kunci Keamanan
  \item \textbf{: Alipay dan WeChat memiliki fitur "Security Lock". Setelah diaktifkan, masuk ke halaman "Dompet" aplikasi memerlukan verifikasi identitas tambahan. Jalur: Alipay Saya Pengaturan Akun \& Keamanan Pusat Keamanan Security Lock.}Pemeriksaan Pembayaran Tanpa Kata Sandi
  \item \textbf{: Banyak orang telah berlangganan berbagai pembayaran tanpa kata sandi, sendiri pun lupa. Ke pengaturan pembayaran Alipay dan WeChat, cari "Pembayaran Tanpa Kata Sandi/Penagihan Otomatis", lihat layanan apa yang telah berlangganan, matikan yang tidak diperlukan.}Manajemen Perangkat Login
  \item \textbf{: Secara berkala periksa perangkat apa saja yang pernah login ke akun Anda. Jika menemukan perangkat yang tidak dikenal, segera keluarkan. Jalur: Alipay Pengaturan Akun \& Keamanan Manajemen Perangkat Login.}4. Bus dan Metro: Ponsel Adalah Kartu Transportasi
\end{itemize}

\booksection{Sekarang naik bus dan metro, sudah jarang orang menggunakan kartu transportasi fisik. Ponsel untuk tap gerbang, ada dua jalur:}

Kartu Transportasi Virtual di Dompet Ponsel

\begin{itemize}
  \item \textbf{: Di Apple Wallet, Huawei Wallet, Xiaomi Wallet langsung buka kartu transportasi virtual. Keuntungan solusi ini: saat tap tidak perlu membuka kunci ponsel, langsung tempel saja, kecepatan respons paling cepat, dan bisa tap tanpa jaringan. Kartu transportasi virtual di sebagian kota juga memiliki diskon transit sama dengan kartu fisik.}Kode Transit Alipay/WeChat
  \item \textbf{: Cari "Transit Code" di aplikasi, aktifkan kode transit bus atau metro kota Anda. Keuntungan solusi ini: tidak perlu top-up, langsung didebit dari saldo atau kartu bank yang terhubung. Namun perlu internet, dan sebelum tap harus buka aplikasi mencari kode.}Saya pribadi lebih merekomendasikan kartu transportasi virtual dompet ponsel. Meski pertama kali perlu top-up 10-20 yuan, tetapi penggunaannya jauh lebih nyaman, terutama saat buru-buru naik metro, tidak perlu berhenti membuka kunci ponsel mencari aplikasi.
\end{itemize}

5. Pembayaran Tagihan, Makanan, Taksi

\booksection{Pembayaran Air/Listrik/Gas}

\begin{itemize}
  \item \textbf{: Di "Life Payment" Alipay atau WeChat, hubungkan nomor akun air, listrik, gas rumah Anda. Setelah itu setiap bulan menerima tagihan, langsung ketuk untuk membayar. Nomor akun biasanya bisa ditemukan di tagihan kertas rumah Anda, atau telepon customer service.}Perbandingan Harga Makanan
  \item \textbf{: Harga dan promo toko yang sama di Meituan dan Ele.me sering berbeda. Sebelum memesan, lihat kedua sisi, kadang bisa hemat beberapa yuan. Selain itu, fitur "Pin Hao Fan" Meituan, memesan bersama teman, harganya bisa jauh lebih murah.}Keamanan Taksi
  \item \textbf{: Di DiDi atau Gaode Taxi, ingat untuk mengatur "Emergency Contact" dan "Trip Sharing". Saat memesan taksi, aplikasi akan otomatis membagikan perjalanan Anda ke kontak yang Anda atur. Selain itu, taksi larut malam prioritaskan "Taksi" bukan "Express", taksi memiliki regulasi industri, tingkat keamanan lebih tinggi.}6. Urus Asuransi Kesehatan, Jaminan Sosial, dan Pemerintahan dengan Ponsel
\end{itemize}

\booksection{Perubahan beberapa tahun ini sangat besar, banyak urusan yang dulu harus antri di kantor, sekarang selesai dalam beberapa menit di ponsel.}

Voucher Elektronik Asuransi Kesehatan

\begin{itemize}
  \item \textbf{: Cari "Medical Insurance Electronic Voucher" di Alipay atau WeChat, setelah diaktifkan, saat ke rumah sakit atau beli obat di apotek langsung tap ponsel, tidak perlu bawa kartu asuransi fisik. Pertama perlu verifikasi pengenalan wajah, setelah itu setiap kali digunakan akan menghasilkan kode dinamis.}Pencarian Jaminan Sosial
  \item \textbf{: Unduh aplikasi "Zhangshang 12333" atau gunakan "Citizen Center" di Alipay, bisa mencari catatan pembayaran jaminan sosial, saldo akun pension, pengajuan tunjangan pengangguran, dll.}Pajak Penghasilan
  \item \textbf{: Rekonsiliasi pajak penghasilan tahunan Maret hingga Juni, bisa diselesaikan di aplikasi "Individual Income Tax". Ingat untuk periksa apakah ada pengurangan khusus tambahan yang bisa diajukan (pendidikan anak, bunga KPR, menafkahkan orang tua, dll.), semua ini bisa membantu Anda mengurangi pajak atau mendapat pengembalian.}Ingat Ini, Cegah Penipuan dan Lindungi Diri
\end{itemize}

\begin{tcbwarn}{Polisi, jaksa, bank, dan customer service tidak akan meminta Anda transfer ke "akun aman".}
\begin{itemize}
  \item Siapa pun yang meminta Anda screenshot kode pembayaran dan kirimkan, adalah penipu.
  \item Jangan klik tautan asing, terutama SMS seperti "ETC kedaluwarsa" "kartu jaminan sosial dinonaktifkan".
  \item Pengembalian dana belanja online akan kembali melalui jalur asal, tidak perlu Anda memberikan nomor kartu bank.
  \item Lakukan Sekarang: Pemeriksaan Keamanan Aplikasi Pembayaran Anda
\end{itemize}

\end{tcbwarn}

\begin{tcbpractice}{Buka Alipay dan WeChat, pastikan pembayaran sidik jari/wajah telah diaktifkan.}
\begin{enumerate}
  \item Temukan fitur "Security Lock" dan aktifkan.
  \item Periksa daftar "Pembayaran Tanpa Kata Sandi/Penagihan Otomatis", matikan layanan yang tidak lagi diperlukan.
  \item Lihat daftar "Perangkat Login", hapus perangkat yang tidak dikenal.
  \item Aktifkan voucher elektronik asuransi kesehatan Anda, sehingga kunjungan rumah sakit berikutnya tidak perlu bawa kartu fisik.
  \item Bab 11: Fotografi dan Manajemen Gambar
\end{enumerate}

\end{tcbpractice}
